\section{Neural-Operator Realization}\label{app:min-neural}

\subsection{Scope}

This appendix explains what the neural-operator layer contributes to the
$(f,g)$ interface. It is a representation family for the maps we want to
learn: build a leaf state, merge child states, and expose a state that the
readout can score at the root. The certificate still comes from the local laws
and the audit. Neural-operator approximation says that the relevant maps can
be represented on the realized calls of a fixed tree; it does not by itself
make a trained checkpoint mergeable.

\subsection{Realized-Call Approximation}

For a fixed tree, collect the actual leaf inputs, merge inputs, and on-range
summary inputs the run makes. If an ideal local-law-satisfying operator is
uniformly approximated on compact neighborhoods of those call sets, transfer
moduli convert approximation tolerance into
$(\varepsilon_{\mathrm{leaf}},\varepsilon_{\mathrm{merge}},
\varepsilon_{\mathrm{idemp}})$ budgets. These budgets feed the approximate
version of the preservation and audit statements.

\subsection{Projection View}

The training objective can be written schematically as
\[
J_\lambda(g)
=
(1-\lambda)L_{\mathrm{oracle}}(g)
+\lambda P_{\mathrm{laws}}(g),
\]
where the projection penalty $P_{\mathrm{laws}}$ vanishes exactly on the
realized local-law subspace. A faithful projection penalty is equivalent
to membership in the exact local-law subspace.

\subsection{Overlap with Classical Sketches}

Classical sketches inhabit the exact local-law subspace when their state,
merge, and readout witnesses are supplied. HLL is the running example: the
ideal state map serializes registers, the merge is pointwise maximum, and the
readout is the cardinality estimator. The scalar cardinality estimates do not
merge; the register state does. Finite attention stacks and FNOs sit on the
learned-operator side. The overlap is the target: learned operators that
behave like valid task-relative state merges on the realized tree.

\subsection{Agarwal Notation}

For compatibility with \citet{Agarwal2013MergeableTODS}, fix
$\varepsilon$ before the merge process. A summary method may output many valid
states, so \(S(D,\varepsilon)\) denotes any valid $\varepsilon$-summary for
the dataset or stream multiset $D$. The size profile is
\[
k(n,\varepsilon)
  =
  \max\{\, |S(D,\varepsilon)| : |D|=n,\ S(D,\varepsilon)\text{ valid}\,\}.
\]
Mergeability means that an algorithm $A$ takes two valid child summaries and
returns a valid summary for the multiset addition of their represented data:
\[
A\big(S(D_1,\varepsilon),S(D_2,\varepsilon)\big)
  \in S(D_1\uplus D_2,\varepsilon),
\qquad
|A(\cdot,\cdot)|\le k(|D_1|+|D_2|,\varepsilon).
\]
If the size bound is independent of $n$, we write it as $k(\varepsilon)$.
The streaming update is the special case \(|D_2|=1\); full mergeability allows
arbitrary child sizes and arbitrary merge trees. The important distinction for
C-TreePO is that \(\uplus\) combines represented data or states, not scalar
query answers. The readout is applied after state merging.

For randomized summaries, validity is an event. After the whole merge tree has
been evaluated, the root state is required to be a valid
\(S(D,\varepsilon)\) with the paper's success probability. C-TreePO preserves
that probability statement: on the root-validity event, the deterministic
readout is correct.

For range-space summaries, the paper's subset notation is literal: an
$\varepsilon$-approximation of \((D,\mathcal{R})\) is a subset
\(S\subseteq D\) such that
\[
\sup_{R\in\mathcal{R}}
\left|
  \frac{|S\cap R|}{|S|}
  -
  \frac{|D\cap R|}{|D|}
\right|
\le \varepsilon .
\]
This subset/readout notation is separate from multiset addition
\(D_1\uplus D_2\), which is the operation represented by the merge tree.

The conversion from C-TreePO to this interface is explicit. Choose the stream
representation, fix $\varepsilon$, choose the state type, expose
\(\mathrm{build}\), \(\mathrm{merge}\), validity \(V_\varepsilon(D,s)\),
state size and \(k(n,\varepsilon)\), then prove build validity, merge closure,
size-profile validity, and valid-state readout correctness. In Lean these are
the fields of \texttt{CTreePOToAgarwalTransform}; its methods produce the
Agarwal state-level summary, the C-TreePO
\path|RelationalMergeablePreferenceShape|, the valid-sized leaf and merge
states, the arbitrary-tree root state, the readout theorem, and the
\(k(|D|,\varepsilon)\) size bound. The randomized version is represented by
\path|RandomizedRelationalMergeablePreferenceShape|.

\subsection{Universal Approximation and Learned Sketches}

The statement that neural operators can capture mergeable sketches is a
capacity statement plus a certification statement. Fix a classical
state-level summary
\[
\mathrm{build}:D\to S,\qquad
\mathrm{merge}:S\times S\to S,\qquad
\mathrm{readout}:S\to Y,
\]
and encode the bounded state space $S$ inside the C-TreePO operator domain.
On the compact set of leaf, merge, and on-range calls realized by a fixed
tree, a neural-operator class with a universal approximation guarantee can
approximate the builder, merger, and readout-facing state map. That is only
the representation step.

The C-TreePO local laws are the certification step. C1 checks that leaf calls
build sufficient states; C3 checks that neural merges preserve the same state
as raw union; C2 checks stability when an already-produced state is reused.
If those losses vanish, the learned operator lies in the exact state-level
mergeable-sketch subspace; in Lean this is the content of
\path|mergeableSketchSummaryClass_subset_exactLocalLawSubspace| and the
bundle theorem \path|local_laws_of_sketch|. If the losses are small, the
root distortion is bounded by the corresponding local-law budget. With a
calibrated readout
$\widehat f$ satisfying
\(\sup_x d(\widehat f(x),\fstar(x))\le \varepsilon_f\), the true-oracle
statement adds the standard \(2\varepsilon_f\) slack. Universal approximation
does not assert that optimization finds this point; the local-law-weighted
objective is the practical projection toward it.
